\newcommand{\figuraPenduloVagon}{%
\begin{figure}[H]
    \centering
    \begin{tikzpicture}[>=Latex,font=\small,line cap=round,line join=round]
        \foreach \y/\titulo in {4.2/{(a) Marco inercial},0/{(b) Marco del vagón}} {
            \draw[very thick,fill=orange!5] (0.4,\y+0.45) rectangle (11.6,\y+3.55);
            \draw[very thick] (0.4,\y+0.45)--(11.6,\y+0.45);
            \foreach \x in {1.5,10.5} {
                \draw[fill=gray!40] (\x,\y+0.25) circle (0.28);
                \draw[fill=white] (\x,\y+0.25) circle (0.10);
            }
            \node[font=\bfseries] at (2.2,\y+3.22) {\titulo};
            \draw[very thick] (6,\y+3.55)--(5.25,\y+1.35);
            \draw[densely dashed] (6,\y+3.55)--(6,\y+1.0);
            \fill[blue!65] (5.25,\y+1.35) circle (0.20);
            \draw[->,ultra thick,blue!70!black] (5.25,\y+1.35)--(5.92,\y+3.3)
                node[midway,left] {$\vec T$};
            \draw[->,ultra thick] (5.25,\y+1.35)--(5.25,\y+0.62)
                node[right] {$m\vec g$};
            \draw (6,\y+2.82) arc[start angle=-90,end angle=-109,radius=0.65];
            \node[fill=orange!5,inner sep=1pt] at (5.73,\y+2.72) {$\theta$};
        }
        \draw[->,ultra thick,green!45!black] (8.2,7.15)--(10.4,7.15)
            node[midway,above] {$\vec a$};
        \draw[->,ultra thick,red!75!black] (5.25,1.35)--(4.0,1.35)
            node[midway,above] {$-m\vec a$};
    \end{tikzpicture}
    \caption{La inclinación del péndulo tiene dos interpretaciones
    equivalentes. En el marco inercial, la tensión produce la aceleración
    horizontal; en el vagón se introduce la fuerza ficticia $-m\vec a$.}
    \label{fig:pendulo-vagon}
\end{figure}%
}